% 第5章第5节：本章小结。

\section{本章小结}
\label{sec:rkhs-summary}

本章从“有限个点值能否作为关于函数的可靠信息”出发，研究点值读取连续的Hilbert函数空间。核与空间之间有双向的对应：空间中每个点的读取有界时，Riesz表示把它写成与核截面$K_x$的内积，核$K(x,z)=\langle K_z,K_x\rangle$对称且半正定；反过来，任何半正定核都唯一地确定一个以它为再生核的空间$H_K$，构造经过系数组合、商去零长度方向、完备化和点值表示四步，核截面的有限组合在其中按核范数稠密。特征映射和核运算把“验证半正定”变成可以操作的检查：写成某个特征向量的内积，或由已知核经非负组合、乘积、换输入、缩放和极限得到。

几项看似相近的性质由不同的条件控制，不能互相代替。固定位置的点值有界由$K(x,x)$给出，全域统一控制要求核的对角有界；读数随位置连续要求核截面距离$\|K_x-K_z\|$随位置趋于零，离散空间$\ell^2(X)$说明前者不蕴含后者。有限观测下，Gram矩阵的零空间只表示系数的冗余，同一函数可以有多种写法；读数相同而函数不同，则来自观测子空间的正交补，它在无限维空间中总存在。投影的唯一性不等于系数的唯一性，最小范数插值由投影几何保证，与系数怎样选取无关。

加入测度后，紧度量空间上的连续半正定核产生紧、自伴、正的积分算子，Mercer定理把核本身展开为一致收敛的谱级数，特征函数在测度全支撑时有唯一的连续代表。核范数在谱坐标下按$1/\lambda_j$加权，于是积分误差小不等于核范数误差小，零谱分量或发散的加权和都把函数排除在$H_K$之外。改变测度会改变特征值和$L^2$坐标，核所定义的内积不变。截断核与截断给定函数控制的是不同的量：前者的误差只依赖核与测度，后者还取决于目标的谱尾；按积分谱截断与按观测中心投影保留的子空间由不同的信息决定，在不同的误差度量下各有优劣。

下一章利用这些空间与表示讨论具体的逼近问题：核截面的有限和能否逼近核空间之外的连续函数，以及怎样用有限个代表统一覆盖整类函数。由核矩阵出发的正则化拟合需要另外给定目标与惩罚，按样本抽取中心的数值近似需要抽样模型，用核比较概率分布则需要把概率测度嵌入核空间；这些问题分别属于优化与模型、统计推断和概率分布的研究，用核比较分布的方法可参考\BookChapterRef{chap:information-theory-and-statistical-geometry}{卷四《数学准备：流形、概率分布与统计推断》“概率分布的分析与优化”章中的最大均值差异}。
